\scanpage{4}
% Typescript face; read after a half turn. The gaps for signs are visible on the sheet.
est toujours vérifiée si $\deleted{\unsure{T}}G$ est un tore). Alors la démonstration donnée montre que dans l'énoncé de 4.6. \inserted{et les corollaires 4.7. et 4.8. ,} on peut se borner à considérer les sous-groupes de la forme $ {}_{m^{r}}\!G$, avec $r\quad 0$. D'ailleurs, on voit aussitôt \deleted{que} \inserted{lorsque $G$ est \uline{simple} sur $S$, alors} pour tout point $s\quad S$, il existe un voisinage ouvert $U$ de $s$ \deleted{tel} et un entier $m\quad 0$, premier à toutes les caractéristiques \deleted{\unsure{rési|rés} \unsure{résiduelle}} résiduelles de $U$, et qui satisfait les conditions précédentes pour $G\quad U$. (Prendre par exemp\deleted{\unsure{o|b}}le pour $m$ l'ordre de la torsion du type de $\deleted{\unsure{H}}G$ \deleted{dans} en $s$). Sous ces conditions, on trouve des $ {}_{m^{r}}\!G$ qui sont finis et \uline{étales} sur $S$, et dont la famille est schématiquement dense, à condition de se restreindre au dessus de $U$. Cette remarque permet dans certains cas (notamment \deleted{dans l'application} \inserted{ceux faisant intervenir} des théorèmes 3.2. et 3.2. bis, \unsure{cf} \inserted{Appendice -} mais non celle des théo\deleted{\unsure{èr}}rèmes 3.6. et 3.6. bis) de se passer d\deleted{\unsure{e}}u \deleted{\unsure{l}}théorè\deleted{\unsure{e}}me 3.1. faisant intervenir la cohomologie de Hochschild. \deleted{Bien entendu, la remarque analogue s'applique au deux corollaires 4.7. et 4.8.}

5. \uline{Homomorphismes centraux des groupes de type multiplicatif}.

\uline{Théorème} 5.1. Soient $u,v:H\quad G$ deux homomorphismes de $S$-préschémas en groupes, avec $H$ de type multiplicatif et de type fini; \deleted{On suppose} $\deleted{\unsure{H}}G$ de présentation finie sur $S$, \deleted{\unsure{ou|on} $S$ localement noethérien\unsure{.}} Soit $s\quad S$ tel que $u_{s}=v_{s}$, et supposons $u$ \uline{central}. Alors il existe un voisinage ouvert $U$ de $s$ tel que $u_{U}=v_{U}$.

Distinguons les deux cas: a) $S$ localement noethérien, auquel cas on peut se borner à supposer $H$ localement de type fini sur $S$. Cette dernière hypothèse et EGA IV 8 nous permettent de remplacer $S$ par le spectre de $\underline{O}_{S,s}$, donc de supposer $S$ local de point fermé $s$. Quitte à remplacer l'anneau local $A$ par son complété, ce qui introduit un changement de base fidèlement plat et quasi-compact, on peut même supposer $A$ complet. En vertu de 3.4. on aura $u_{n}=v_{n}$ pour tout $n$,

% archive sheet number: 3

